% Generated draft: requires mathematical and source review.
\documentclass[11pt]{article}
\usepackage[margin=25mm]{geometry}
\usepackage{fontspec}
\setmainfont{Latin Modern Roman}
\usepackage{amsmath,amssymb,mathtools,unicode-math}
\setmathfont{Latin Modern Math}
\usepackage{enumitem,needspace,xurl,hyperref}
\setlength{\parindent}{0pt}
\setlength{\parskip}{5pt}
\setlength{\emergencystretch}{3em}
\newcommand{\N}{\mathbb{N}}
\newcommand{\Z}{\mathbb{Z}}
\newcommand{\Q}{\mathbb{Q}}
\newcommand{\R}{\mathbb{R}}
\newcommand{\C}{\mathbb{C}}
\newcommand{\F}{\mathbb{F}}
\newcommand{\A}{\mathbb{A}}
\newcommand{\PP}{\mathbb{P}}
\newcommand{\E}{\mathbb{E}}
\newcommand{\eps}{\varepsilon}
\newcommand{\dd}{\,\mathrm{d}}
\begin{document}
\section*{UnsolvedMath ID: 30003844}
\textbf{Source problem code:} catalogue entry
\subsection{Definitions and mathematical statement}
The supplied record concerns subgroups of \(\operatorname{PL}_o(I)\). Under the usual interval convention, \(I=[0,1]\) and \(\operatorname{PL}_o(I)\) is the group, under composition, of increasing piecewise-linear homeomorphisms of \(I\) with finitely many breakpoints. These interval and breakpoint conventions are not explicitly specified in the supplied record and require confirmation from the source. In this model, Thompson's group \(F\) consists of the maps with dyadic rational breakpoints and slopes that are integral powers of \(2\). A group is \emph{\(F\)-less} if it has no subgroup isomorphic to \(F\).

Two groups \(G\) and \(K\) are \emph{biembeddable}, written here as \(G\sim K\), if there exist injective group homomorphisms \(G\to K\) and \(K\to G\). For a subset \(A\) of a group, \(\langle A\rangle\) denotes the subgroup it generates.

The supplied record does not define a \emph{fast generating set}. Consequently, the following is a statement schema with an unresolved predicate, rather than a fully specified conjecture. Interpreting ``a finite fast generating set \(A\) with the same biembeddability class'' as referring to the group \(\langle A\rangle\), the first question is whether
\[
\begin{gathered}
\text{for every finitely generated }H\leq\operatorname{PL}_o(I)
\text{ containing no subgroup isomorphic to }F,\\
\text{there exists a finite fast set }A\subseteq\operatorname{PL}_o(I)
\text{ such that }H\sim\langle A\rangle.
\end{gathered}
\]
Here ``fast'' must be supplied with its intended technical definition before this schema becomes a precise mathematical statement. The wording does not resolve whether \(A\) must instead be a generating set of \(H\) itself. Under that reading one would require \(A\subseteq H\) and \(\langle A\rangle=H\), and the biembeddability condition would be automatic. These readings must not be identified without source confirmation.

The second question asks for a useful characterization of fast subgroups \(K\leq\operatorname{PL}_o(I)\) containing no subgroup isomorphic to \(F\). The record supplies neither the definition of a fast subgroup nor a required form of characterization. It also does not explicitly restrict this second question to finitely generated subgroups. These qualifications remain unresolved.
\subsection{Short English statement}
The problem asks whether every finitely generated subgroup of \(\operatorname{PL}_o(I)\) containing no copy of Thompson's group \(F\) is biembeddable with a group having a finite fast generating set, and seeks a characterization of fast subgroups containing no copy of \(F\). The technical meaning of ``fast'' and whether the generating set must generate the original subgroup require clarification from the source.
\subsection{Sources}
\begin{itemize}
\item Local UnsolvedMath registry, permanent ID \(30003844\), title ``Fast Generators for F-Less Homeomorphism Groups,'' exact source problem code \texttt{catalogue entry}, in the Oberwolfach Reports Open Problems collection. The record gives proposed year 2018 and identifies no individual proposer.
\item Supplied report citation: \emph{Cohomological and Metric Properties of Groups of Homeomorphisms of R} (2019), Oberwolfach Reports, pp.\ 1579--1633, DOI \texttt{10.4171/owr/2018/26}. \url{https://doi.org/10.4171/owr/2018/26}. The background describes these reports as records of mathematical workshops and their open-problem discussions.
\item Supplied literature reference labeled ``Brin'': \url{https://arxiv.org/abs/math/0501085}.
\item Supplied literature reference labeled ``Bleak--Brin--Moore'': \url{https://doi.org/10.4171/JCA/49}.
\item Supplied literature reference labeled ``Oberwolfach Report 2019'': \url{https://ems.press/content/serial-article-files/46748}.
\item Literature assessment dated exactly \texttt{2026-08-21}: the supplied curation labels the problem open, describes the Brin--Sapir conjecture as remaining open in general, and reports partial structure from restricted analyses of elementary amenable subgroups and fast generating sets, without a proof. This is supplied dated curation derived from the OWR record, not an independently verified literature assessment. The record itself requests independent verification of time-sensitive status and cited claims.
\item The registry labels the statement \texttt{corrected\_verified} and reports a manual semantic, typographic, and LaTeX review asserting self-containment and preservation of the target. These are supplied review claims; the unresolved terminology and interpretation identified here still require mathematical source review.
\end{itemize}
\end{document}
